\section{Uscita, profilo di chi porta il sensore, perdita di affidabilità, condizioni di validità e limiti}\label{sec:lim}

Questo capitolo descrive ciò che il sistema dichiara verso l'esterno e ciò che non dichiara mai, come usa la configurazione di chi porta il sensore, come segnala che non è più affidabile, e dentro quali condizioni valgono gli enunciati dei capitoli precedenti. Nessuna uscita è un comando: il sistema osserva, sorveglia lo stato della scena e avvisa; non ha funzioni d'ingaggio né di puntamento. Tutte le cifre sono dimostrazioni in un modello dichiarato, calcoli o simulazioni a sensore fermo: nessuna è una misura su un sensore fisico.

\subsection{Contratto d'uscita}\label{sec:lim-uscita}

\paragraph{Che cosa si dichiara.} Di ogni cella della mappa polare il sistema dichiara uno fra quattro stati: \emph{libero}, \emph{occupato osservato}, \emph{in ombra}, \emph{non osservato} (\S\ref{sec:map}); «non osservato» non diventa mai «libero». Di ogni oggetto seguito dichiara se lo vede, se non lo vede ma indica dove può essere (con un livello, quando è dichiarabile), o se non lo sa. Gli stati di una traccia e i loro segni sono nella tabella~\ref{tab:lim-stati} e nella figura~\ref{fig:lim-segni}. Valgono due regole di forma.
\begin{enumerate}[label=(\alph*)]
\item \emph{Contorno pieno solo con una misura.} Il disco pieno compare soltanto nei cicli in cui la traccia ha una misura non tagliata dentro la porta; una previsione non esce mai con il segno di un'osservazione, e il passaggio avviene in un solo verso (da previsto a osservato serve una misura).
\item \emph{Regione tratteggiata con l'età del dato.} Ogni regione di previsione porta l'etichetta «previsto · livello · larga $X{,}X\unita{m}$ · non visto da $Y{,}Y\unita{s}$»; l'età è misurata dall'ultima misura non tagliata e non è mai scritta come nulla.
\end{enumerate}
Ogni distinzione sta nella forma, nella trama e nella parola; il colore la ripete soltanto e non è mai verde. Ogni riga d'uscita porta contatore del fotogramma, età, livello, configurazione e l'indicazione «osservato» o «dedotto».

\begin{table}[H]\centering\small
\caption{Stati di una traccia e uscita corrispondente. Le soglie sono in secondi; a $15\unita{Hz}$ tre fotogrammi valgono $0{,}2\unita{s}$.}\label{tab:lim-stati}
\begin{tabularx}{\linewidth}{@{}p{2.5cm}Lp{4.6cm}L@{}}
\toprule
Stato & Quando & Segno e parole & Finisce quando\\
\midrule
osservato & misura non tagliata in porta nel ciclo & disco pieno; «visto · $d\unita{m}$» & manca la misura\\
visto in parte & il gruppo è tagliato da un bordo d'ombra o dal bordo del campo & contorno pieno, senza velocità, insieme a regione e avviso & torna intero o scompare\\
non confermato & traccia confermata, senza misura, in campo aperto & anello vuoto con l'età, per $0{,}2\unita{s}$ & evento «non più rilevato»\\
previsto & entrata in un'ombra, livello dichiarabile & anello vuoto, regione tratteggiata $\Reg_\alpha(t)$, età & diventa «previsto ampio»; non torna indietro\\
previsto ampio & regione oltre l'$80\,\%$ dell'ombra, oltre l'orizzonte del livello, ingresso da fermo, sensore non in salute & tutta la parte dell'ombra non vista libera, a tratteggio rado, senza punto & riaggancio; ombra vista libera; uscita dall'area osservata\\
possibile presenza & una traccia è entrata nell'ombra & ombra più fascia $w$, quadrettata, con l'età; «residua» dopo un riaggancio & solo dove il suolo è rivisto libero oltre $w$ dall'ombra corrente\\
\bottomrule
\end{tabularx}
\end{table}

\begin{figure}[H]\centering
\begin{tikzpicture}[x=1cm,y=1cm,font=\footnotesize]
\fill (0,0) circle (0.16); \node[anchor=west] at (0.3,0) {osservato};
\draw[line width=0.9pt] (2.55,-0.18) rectangle (2.95,0.18); \node[anchor=west] at (3.1,0) {visto in parte};
\draw[dashed,pattern=north east lines,pattern color=muted] (6.0,-0.3) ellipse (0.55 and 0.32);
\draw[line width=0.8pt,fill=white] (6.0,-0.3) circle (0.13); \node[anchor=west] at (6.7,-0.05) {previsto: anello vuoto,};
\node[anchor=west] at (6.7,-0.42) {regione tratteggiata, età};
\draw[dotted,line width=0.8pt,pattern=north west lines,pattern color=muted!60] (10.6,-0.6) -- (11.9,-0.6) -- (11.6,0.15) -- (10.9,0.15) -- cycle;
\node[anchor=west] at (12.0,-0.05) {previsto ampio:};
\node[anchor=west] at (12.0,-0.42) {nessun punto};
\draw[line width=1pt,pattern=crosshatch,pattern color=muted] (0,-1.55) -- (1.6,-1.55) -- (1.25,-0.75) -- (0.35,-0.75) -- cycle;
\node[anchor=west] at (1.8,-0.97) {possibile presenza: ombra dell'ostacolo più la fascia $w$, con l'età;};
\node[anchor=west] at (1.8,-1.34) {non dipende dall'identità e finisce solo dove il suolo è rivisto libero};
\end{tikzpicture}
\caption{Segni dell'uscita visiva. Il segno pieno è riservato a ciò che ha una misura nel ciclo corrente.}\label{fig:lim-segni}
\end{figure}

\paragraph{Parole per settore.} La zona valutata è divisa in 17 settori d'azimut; ogni settore riceve un solo livello, il più alto fra: «attenzione» (un elemento mobile in avvicinamento, detto per primo, poi un ostacolo occupato a $d$ metri); «zona non verificata» (possibile presenza, regione di previsione, ombra, o un buco di dati prima della portata); «suolo visto libero da $1{,}3\unita{m}$ a $d\unita{m}$», dove $d$ è il raggio della prima cella non libera del settore e non supera la portata meno la fascia, $3-0{,}40=2{,}6\unita{m}$; «sconosciuto» se il settore non vede il suolo da $1{,}3\unita{m}$. Sopra tutti, fuori scala, sta «affidabilità persa» (\S\ref{sec:lim-salute}). Con la geometria del capitolo~\ref{sec:geo}, nove settori su 17 vedono il suolo da $1{,}3\unita{m}$; gli 8 laterali dicono «sconosciuto» per geometria \prova{S}. La frase per settori è una regola di progetto: non è stata generata in alcuna prova, e l'ordine fra i primi due livelli non è stato provato \prova{N}.

\paragraph{Che cosa non si dichiara mai.} (i) «Via libera» o «nessun ostacolo osservato fino a $d$»: la frase parla del suolo letto, non del volume, perché un oggetto basso vicino, una barra sottile o un ostacolo sospeso non interrompono la lettura del suolo (\S\ref{sec:lim-limiti}). (ii) Una sosta, un'uscita o un'identità certa: dietro un ostacolo non si distinguono; l'identità al ritorno è «presunta» o «ambigua». (iii) Un livello numerico oltre $3\unita{s}$ senza osservazione, con il predittore di ripiego, per chi entra nell'ombra da fermo o con il bordo dell'ombra non riconosciuto. (iv) Massa o accelerazione stimate (\S\ref{sec:sti}). (v) Una confidenza in percentuale o un intervallo «$\pm$ metri» attorno al punto: il punto previsto non ha una precisione dichiarata, la ha la regione. (vi) Una classe senza la sua origine: nelle prove l'etichetta è «assegnata a mano» e una riga fissa dichiara che nessun modello di classificazione è in esecuzione.

\paragraph{Due profili d'uscita.} Le regole sono le stesse; cambia il canale. Nel profilo di \emph{monitoraggio di un varco o di un'area di lavoro} l'uscita è lo schermo con il registro degli eventi. Nel profilo di \emph{avviso di presenza per una persona con disabilità visiva ferma} accanto al sensore l'uscita è sonora: un segnale di $0{,}3\unita{s}$ per ogni cambio di stato e la frase «direzione, distanza, stato» solo a stato stabile, con precedenza (1) affidabilità persa, ripetuta finché dura, (2) elemento visto in avvicinamento, (3) scomparso dietro l'ostacolo, (4) previsto ampio e poi non ricomparso, (5) ricomparso, (6) identità ambigua, (7) la frase. All'avvio il canale dichiara la sorgente dei dati e che il silenzio non significa assenza. Il canale sonoro è una specifica: che i segnali si distinguano, che diano una posizione e con quale ritardo arrivino non è stato provato con persone \prova{N}.

\subsection{Registro degli eventi}\label{sec:lim-registro}

Per il monitoraggio di attività il sistema scrive una riga per evento: ora, identità, evento, \emph{come lo sa} («osservato» oppure «dedotto»), età del dato, settore, configurazione. Gli eventi sono: ingresso; scomparsa dietro l'ostacolo; «non ricomparso da $n\unita{s}$» (dedotto); ricomparsa con identità presunta o ambigua; affidabilità persa e tornata; con le classi, etichetta assegnata, cambiata, smentita dalla cinematica, classe fuori scenario. Nessun evento afferma una sosta o un'uscita, e nessun evento nasce da un peso del predittore. La logica del registro è stata controllata in simulazione su scene sintetiche con una catena ridotta: 0 violazioni delle regole del registro su 600 scene; eventi mancati nel 16--52\,\% delle scene e scomparse spurie nel $12{,}5\,\%$ \prova{S}. Non è stata eseguita sulla catena completa \prova{N}.

\subsection{Salute del sensore e perdita di affidabilità}\label{sec:lim-salute}

\paragraph{Controllo all'avvio.} A scena vuota e ferma: (1) i filtri della camera che ricordano o riempiono i dati mancanti (persistenza temporale, riempimento dei buchi \cite{realsense-filtri-post}) sono spenti; (2) autoverifica: coprendo il sensore i pixel validi devono scendere sotto il $10\,\%$ entro 3 fotogrammi; (3) la posa viene dalla taratura, non è mai stimata dai primi fotogrammi (su una registrazione pubblica \cite{tum-rgbd} quella stima sbaglia di $5^\circ$ \prova{P}) e si confronta con una misura diretta: uno scarto di $0{,}5^\circ$ o più impedisce l'avvio; (4) un oggetto alto $10\unita{cm}$ a $2{,}5\unita{m}$ deve risultare «occupato»; (5) sul suolo vuoto, in $30\unita{s}$, non devono nascere gruppi né tracce confermate; la prova si esegue con i gruppi a celle forti (\S\ref{sec:sti-mobile}): con quelli della configurazione di riferimento, su profondità simulata dalle schede tecniche, nascono 137 e 206 tracce confermate in 9 sequenze e la prova non sarebbe superata \prova{S}. Se cade una delle prime quattro prove il sistema non parte; se cade la quinta restano spenti i gruppi a celle forti, le condizioni aggiuntive del riaggancio, l'avviso senza traccia e l'identità al ritorno, e l'uscita lo dichiara. Il motivo della (1): un filtro che riempie i buchi trasforma un dato mancante in un dato valido, cioè in un possibile «libero»; su profondità simulata dalle schede tecniche, con i filtri accesi e il sensore coperto, restano «libero» 989--995 celle per altri 3 fotogrammi \prova{S}. Il motivo della (3) e della (4) è nel limite~\ref{lim:posa}.

\paragraph{Indicatori in esercizio.} Il sensore è dichiarato non in salute se: non arriva un fotogramma nuovo per $0{,}13\unita{s}$ o ne arrivano tre identici; i pixel validi restano sotto la metà del valore di taratura per $0{,}2\unita{s}$; l'unità inerziale segnala più di $0{,}5^\circ$ o di $5^\circ/\mathrm{s}$; il piano del suolo stimato da fermo si scosta di $0{,}5^\circ$ dalla taratura. La quota dei validi si conta sui pixel e non sulle celle: con il $10\,\%$ di pixel validi 1\,444 celle su 1\,460 risultavano ancora note \prova{S}.

\paragraph{Che cosa accade.} «Affidabilità persa» è uno stato del sistema, non di una traccia. La mappa diventa tutta «non osservato»; ogni traccia confermata passa a «previsto ampio · dedotto, sensore non affidabile» e continua a uscire; il ciclo emette comunque, con il segnale dedicato. Al ritorno dopo un tempo $T$ le regioni di previsione e l'area d'avviso sono dilatate di $\vmax T+w$ e intersecate con ciò che non è visto libero nel fotogramma; una traccia la cui intersezione è vuota si chiude come «vista libera», e ogni riaggancio è marcato «identità ambigua». La scelta è quella del guasto verso lo stato più restrittivo, come nel segnalamento ferroviario \cite{us-49cfr-236-205}.

\emph{Livello di prova.} Segnale entro $0{,}13\unita{s}$ e 0 falsi segnali in $60\unita{s}$, in simulazione, con rumore dell'unità inerziale ipotizzato \prova{S}. Il ramo del ritorno è verificato per calcolo in pianta: intersecando con «non osservato» la persona passata dietro l'ostacolo era fuori dalla regione in 13 istanti su 13, intersecando con «non visto libero» era dentro in 13 su 13 \prova{C}; non è stato riprovato sulla catena \prova{N}.

\subsection{Profilo di chi porta il sensore}\label{sec:lim-portatore}

\paragraph{Configurazione.} Chi porta il sensore non è classificato: è descritto da una configurazione deterministica, fissata all'avvio e immutabile in esercizio. I campi sono: identificativo; velocità propria tipica e massima $v_e$, $v_{e,\max}$ (m/s); tempo di reazione $T_r$ (s), dall'avviso emesso alla reazione iniziata, più l'eventuale ritardo di collegamento; decelerazione di frenata $a_f$ e accelerazione massima (m/s$^2$); gradi di libertà e olonomia; canale d'avviso; indicatore «sensore in moto»; elenco delle classi ammesse nel sito (lo \emph{scenario}, \S\ref{sec:sti}).

\paragraph{Formule.} Sia $t_{\mathrm{conf}}=0{,}267\unita{s}$ il tempo di conferma di una traccia (\S\ref{sec:sti}), $\tau=0{,}155\unita{s}$ la latenza stimata (\S\ref{sec:lim-latenza}) e $C=0{,}95\unita{m}$ un margine costante. Per un oggetto che si avvicina frontalmente a velocità $v$ senza frenare, il tempo d'anticipo necessario e la distanza a cui l'oggetto deve essere rilevato sono
\begin{align}
T_{\mathrm{ant}} &= t_{\mathrm{conf}}+\tau+T_r+\frac{v_e}{a_f}, \label{eq:lim-tant}\\
d_{\mathrm{serve}}(v) &= (v+v_e)\,(t_{\mathrm{conf}}+\tau+T_r)+\frac{v_e^2}{2a_f}+\frac{v\,v_e}{a_f}+C. \label{eq:lim-dserve}
\end{align}
Il primo termine è lo spazio di chiusura fino all'inizio della reazione, il secondo lo spazio d'arresto del portatore, il terzo lo spazio percorso dall'oggetto durante la frenata. La forma «velocità per tempo più costante» è quella della distanza di protezione nella sorveglianza di velocità e separazione; $C$ somma $0{,}85\unita{m}$, valore che la stessa fonte riporta dalla norma sul posizionamento dei ripari per l'avvicinamento perpendicolare con capacità di rilevamento fra 40 e $70\unita{mm}$ \cite{marvel2017} (la norma non è stata consultata nel testo originale, e che la capacità di rilevamento del sistema stia in quell'intervallo non è verificato), e $0{,}10\unita{m}$ d'incertezza del sensore: il confronto dà un ordine di grandezza e non è una dichiarazione di conformità. La \emph{velocità coperta} è la $v$ per cui $d_{\mathrm{serve}}(v)$ uguaglia la portata; a sensore fermo ($v_e=0$) e con $3\unita{m}$ di portata vale $(3-C)/(t_{\mathrm{conf}}+\tau+T_r)$.

\paragraph{Uso.} All'avvio il sistema calcola $T_{\mathrm{ant}}$ e $d_{\mathrm{serve}}(\vmax)$ per ogni classe dello scenario; se $d_{\mathrm{serve}}(\vmax)$ supera la portata dichiara in una riga fissa «avviso tempestivo non dichiarato per» quelle classi; se la portata divisa per $\vmax$ è minore di $t_{\mathrm{conf}}+\tau$ dichiara che per quelle classi l'avviso esce dopo il contatto; se $v_e>0$ dichiara «soglie non provate a sensore in moto». In esercizio, a ogni «attenzione» si calcola il tempo disponibile come distanza divisa per la velocità d'avvicinamento osservata aumentata di $2\sigma_v$: se è minore di $T_{\mathrm{ant}}$ si aggiunge la parola «tardivo». Il profilo non abbassa mai un livello e nessuna etichetta di classe è argomento di queste regole.

\begin{table}[H]\centering\small
\caption{Profili d'esempio e grandezze derivate dalle \eqref{eq:lim-tant}--\eqref{eq:lim-dserve}, portata $3\unita{m}$ \prova{C}. Origine dei valori: (i) ipotesi di progetto; (r) valore riportato da una fonte non verificata. Le ultime due colonne sono analisi fuori dal dominio provato \prova{N}.}\label{tab:lim-portatore}
\begin{tabularx}{\linewidth}{@{}Lp{1.9cm}p{1.0cm}p{2.2cm}p{1.7cm}p{2.5cm}@{}}
\toprule
Profilo & $T_r$ (s) & $T_{\mathrm{ant}}$ (s) & Velocità coperta, sensore fermo (m/s) & In moto: $v_e$; $a_f$ & $d_{\mathrm{serve}}(0)$; velocità coperta\\
\midrule
stazione fissa, avviso a una persona vicina & 1,0 (i) & 1,42 & 1,44 & --- & ---\\
veicolo terrestre senza equipaggio, arresto automatico & 0,3 (i) & 0,72 & 2,84 & 1,0 (i); 1,5 (i) & $2{,}01\unita{m}$; $0{,}72\unita{m/s}$\\
lo stesso, con operatore remoto & 1,5 (r) + 0,3 (i) & 2,22 & 0,92 & --- & ---\\
multirotore senza equipaggio in volo stazionario & 0,3 (i) & 0,72 & 2,84 (l'oscillazione non è provata) & 3,0 (r); 3,0 (i) & $4{,}62\unita{m}$: oltre la portata\\
persona con disabilità visiva & 1,5 (i), campo 0,5--1,5 & 1,92 & 1,07 & 1,0 (i); 1,0 (i) & $3{,}37\unita{m}$: oltre la portata\\
\bottomrule
\end{tabularx}
\end{table}

\paragraph{Lettura e limiti.} (a) A sensore fermo del profilo conta solo $T_r$. Con $T_r=1{,}0\unita{s}$ la velocità coperta dall'avviso tempestivo è $1{,}44\unita{m/s}$: una persona a passo normale ($1{,}4\unita{m/s}$) vi rientra con $6\unita{cm}$ di margine sulla distanza, il passo svelto no. Non è la velocità coperta dalla fascia di guardia ($1{,}13\unita{m/s}$, \S\ref{sec:lim-validita}), che il passo normale supera: a $1{,}4\unita{m/s}$ il contenimento nella fascia non è dimostrato. (b) Il tempo di reazione di una persona con disabilità visiva a un avviso sonoro non ha una fonte: $1{,}0\unita{s}$ e $1{,}5\unita{s}$ sono due ipotesi di progetto, e con la seconda la velocità coperta scende a $1{,}07\unita{m/s}$, sotto il passo normale; il valore si decide solo con una misura su persone \prova{N}. Il valore di $1{,}5\unita{s}$ per l'operatore remoto è la reazione di frenata a sorpresa di un conducente riportata in letteratura, non verificata sulla fonte. (c) \emph{Ogni cifra con $v_e>0$ è un calcolo fuori dal dominio provato}: con il sensore in moto l'ombra si sposta, cade l'ipotesi di ostacolo fermo ed entra l'errore di posa (\S\ref{sec:lim-moto}); nessuna soglia del sistema è stata provata in quelle condizioni, nemmeno in simulazione. La fascia di guardia diventerebbe $w=\varepsilon_\Omega+(v_{\mathrm{scen}}+v_e)\cdot0{,}355\unita{s}$, con $\varepsilon_\Omega$ l'errore del bordo d'ombra e $v_{\mathrm{scen}}$ la velocità massima dello scenario. Chi si muove paga il proprio arresto: a $1{,}0\unita{m/s}$ con $a_f=1{,}0\unita{m/s^2}$ e $T_r=1{,}5\unita{s}$ la sola distanza per un ostacolo fermo ($3{,}37\unita{m}$) supera la portata. (d) Le classi veloci sono fuori portata comunque: la tabella del dominio per classe è nel capitolo~\ref{sec:occ}.

\subsection{Bilancio di latenza}\label{sec:lim-latenza}

\begin{table}[H]\centering\small
\caption{Voci note della latenza $\tau$, dall'evento alla scheda audio, a 30 fotogrammi al secondo \prova{C}. Nessuna voce è misurata sulla catena del prototipo \prova{N}.}\label{tab:lim-latenza}
\begin{tabularx}{\linewidth}{@{}p{5.2cm}p{2.0cm}L@{}}
\toprule
Tratto & Valore medio & Natura del dato\\
\midrule
evento $\to$ inizio dell'esposizione & $16{,}7\unita{ms}$ & calcolo: uniforme fra 0 e $33{,}3\unita{ms}$\\
esposizione $\to$ fotogramma disponibile & $95\unita{ms}$ & misura pubblicata dal costruttore sull'intero banco di prova, dal movimento allo schermo, nel modo $848\times100$ a 30 fotogrammi al secondo \cite{realsense-alta-velocita-d435}: stima per eccesso del solo sensore, non misurata alla risoluzione usata qui; media, senza percentili\\
fotogramma $\to$ comando d'uscita & $22\unita{ms}$ & misura su scene sintetiche con codice non ottimizzato\\
campioni $\to$ scheda audio & $21{,}3\unita{ms}$ & un quanto di 1\,024 campioni a $48\unita{kHz}$\\
\midrule
somma delle medie & $155\unita{ms}$ & \\
\bottomrule
\end{tabularx}
\end{table}

La somma si legge in due modi. Come \emph{media} vale $155\unita{ms}$. Come \emph{percentile}: la sola prima voce la fa variare fra 138 e $172\unita{ms}$, e se le altre fossero costanti il 95° percentile sarebbe $170\unita{ms}$; un requisito «95° percentile non oltre $150\unita{ms}$» non è quindi soddisfatto dalle sole voci note. Un secondo dato del costruttore per il tratto del sensore, un unico esempio di circa $71\unita{ms}$, porterebbe la somma a circa $131\unita{ms}$: le due cifre vengono da fonti diverse e nessuna decide. Mancano, perché senza dato, la sintesi vocale (nulla con suoni preregistrati) e il tratto dalla scheda al suono. La latenza entra nella fascia di guardia (\S\ref{sec:occ}) e nelle \eqref{eq:lim-tant}--\eqref{eq:lim-dserve}: ogni cifra che la usa è condizionata alla stima.

\subsection{Invarianti del sistema}\label{sec:lim-invarianti}

\begin{table}[H]\centering\footnotesize
\caption{Invarianti, livello di prova e capitolo che li tratta. «Per costruzione»: la proprietà segue dalla forma della regola e si controlla con un'asserzione a ogni ciclo; non è una misura.}\label{tab:lim-invarianti}
\begin{tabularx}{\linewidth}{@{}rLLp{1.1cm}@{}}
\toprule
 & Enunciato & Livello di prova & Cap.\\
\midrule
1 & «Libero» esiste solo nel fotogramma corrente, dove metà dell'impronta della cella e delle otto vicine legge il suolo; non si ricorda e non si deduce & \prova{S} 0 celle libere su un ostacolo in 300 condizioni (9\,000 fotogrammi), modello ideale del sensore, posa entro $0{,}5^\circ$ (limite superiore per condizione $1{,}3\,\%$); \prova{N} su sensore & \ref{sec:map}\\
2 & Fuori campo, oltre portata, suolo non visibile vicino, ombra, impronta letta per meno di metà (pixel non validi compresi), sensore non in salute: mai «libero» & \prova{D} per costruzione, senza ipotesi sul sensore. Vale per ciò che non è visto; non afferma che ciò che è «libero» sia libero & \ref{sec:map}\\
3 & Segno pieno solo con una misura non tagliata in porta; una previsione non esce come osservazione & per costruzione; soglia 0 non superata: 8 passaggi su 90 e, in una seconda serie, riagganci a soggetto ancora nascosto in 11 e 6 passaggi su 140, con la configurazione di riferimento su profondità simulata dalle schede tecniche; 1 e 1 su 140 con i gruppi a celle forti \prova{S} & \ref{sec:sti}\\
4 & $\Reg_\alpha(t)$ è contenuta nell'area d'avviso e non copre suolo «libero» oltre $0{,}20\unita{m}$ dall'ombra & \prova{D} l'annidamento; \prova{S} la fascia di $0{,}20\unita{m}$ (0 casi misurati, non garantiti dalla sola regola del peso) & \ref{sec:occ}\\
5 & Nessun timer, modello appreso, etichetta di classe, peso, identità o conferma dell'operatore produce un «libero», toglie o restringe un avviso, abbassa un livello & per costruzione; per l'etichetta \prova{S} sulla sola fascia: 0 violazioni su 20\,000 sequenze di etichette e 0 su altre 20\,000; \prova{N} sulla catena & \ref{sec:occ}\\
6 & L'evidenza negativa vale dove il suolo è letto ora; restringe $\Reg_\alpha$, mai l'avviso & per costruzione & \ref{sec:occ}\\
7 & L'avviso finisce solo guardando, cella per cella, oltre $w$ dall'ombra corrente & per costruzione; \prova{S} 0 fotogrammi con il nascosto fuori dall'avviso su 336, 169 e 2\,360, su profondità simulata dalle schede tecniche; che l'avviso termini quando il suolo è rivisto dopo la rimozione dell'ostacolo non è sostenuto: resta acceso dopo $2\unita{s}$ in 10 prove su 10 (72--100\,\%) & \ref{sec:occ}\\
8 & Ogni uscita porta contatore, età, livello, configurazione e «osservato» o «dedotto» & per costruzione & \ref{sec:lim}\\
9 & Massa e accelerazione non sono stimate & \prova{D} nel modello del punto materiale & \ref{sec:sti}\\
10 & Continuità: per ogni traccia con 4 fotogrammi di misura entrata in un'ombra, a ogni ciclo fino alla chiusura escono regione, età e stato & per costruzione, data la premessa; \prova{S} 0 passaggi con un vuoto dopo l'ingresso su 23\,402 (6\,000 camminatori distinti: meno di 7 su 10\,000), modello ideale del sensore; la premessa manca in 928 passaggi, e con l'avviso senza traccia 24 restano senza avviso & \ref{sec:occ}, \ref{sec:val}\\
11 & Configurazione, scenario, profilo del portatore e interruttori sono fissati all'avvio & regola di progetto & \ref{sec:lim}\\
12 & Il livello più basso è «suolo visto libero da $1{,}3\unita{m}$ a $d\unita{m}$», $d\le2{,}6\unita{m}$; altrimenti «sconosciuto» & regola di progetto & \ref{sec:lim}\\
13 & Senza autoverifica, posa e prova della posa il sistema non parte; non in salute, il ciclo esce comunque & \prova{N} & \ref{sec:lim}\\
14 & Ogni evento dice se è osservato o dedotto; nessuno afferma sosta, uscita o identità certa & \prova{S} su 600 scene sintetiche, catena ridotta & \ref{sec:lim}\\
15 & L'etichetta di classe può solo allargare la fascia; senza etichetta, o con etichetta smentita dalla cinematica, vale il profilo prudente dello scenario, senza livello numerico & per costruzione; \prova{N} sulla catena & \ref{sec:sti}, \ref{sec:occ}\\
\bottomrule
\end{tabularx}
\end{table}

Tre forme più forti di questi enunciati non reggono, e per questo non sono fra gli invarianti: che il «libero» ricordato per tre fotogrammi sposti l'oggetto di meno della fascia (con la latenza lo spostamento è $0{,}72\unita{m}$ a $2{,}5\unita{m/s}$ \prova{C}); che la previsione non influenzi l'avviso (la regola d'ingresso nell'ombra usa la previsione: vale solo la monotonia, cioè la previsione accende e non restringe); che «il 100\,\% degli istanti dentro l'avviso» sia una misura (è una conseguenza della costruzione finché l'errore dell'ombra non supera $w$ più mezzo corpo).

\subsection{Condizioni di validità}\label{sec:lim-validita}

\begin{riquadro}{Condizioni di validità degli enunciati e delle cifre}
\begin{enumerate}[label=\arabic*.,itemsep=1pt]
\item \textbf{Camera ferma, posa nota.} Quota $h=1{,}2\unita{m}$, inclinazione $\theta=15^\circ$, errore d'inclinazione minore di $0{,}5^\circ$; suolo piano.
\item \textbf{Al chiuso, fino a $3\unita{m}$.} Oltre la portata non c'è mappa né avviso; la frase del suolo arriva al più a $2{,}6\unita{m}$.
\item \textbf{Superfici che rimandano segnale.} Ostacoli opachi che restituiscono una profondità valida e corretta; filtri della camera spenti.
\item \textbf{Un soggetto per ombra,} visto per almeno 4 fotogrammi prima di entrarvi; ostacolo fermo.
\item \textbf{Velocità entro il limite coperto dalla fascia.} Con $w=0{,}40\unita{m}$, tempo di rilevamento assunto $t_a=0{,}133\unita{s}$, latenza stimata $0{,}155\unita{s}$ e tenuta di un ciclo $0{,}067\unita{s}$, per il centro del soggetto: $0{,}40/0{,}355=1{,}13\unita{m/s}$; $0{,}56\unita{m/s}$ se il bordo dell'ombra sbaglia di $0{,}20\unita{m}$ \prova{C}.
\item \textbf{Ostacolo che nasconde per intero.} Per una persona alta $1{,}70\unita{m}$ a $2{,}3\unita{m}$ dal sensore servono da $1{,}57$ a $1{,}61\unita{m}$ d'ostacolo con il pannello fra $1{,}7$ e $1{,}9\unita{m}$ (Tabella~\ref{tab:geo-amin}).
\item \textbf{Per il livello numerico della regione:} ostacolo largo almeno $2\unita{m}$, persona che resta sul proprio percorso, bordo dell'ombra noto esattamente, non oltre $3\unita{s}$, classe «persona».
\end{enumerate}
Fuori da queste condizioni ogni cifra perde il suo dominio; restano le proprietà che valgono per costruzione sul singolo fotogramma: ciò che non è osservato non è mai «libero», e nessuna memoria, etichetta o timer produce un «libero».
\end{riquadro}

\subsection{Limiti}\label{sec:lim-limiti}

Ogni limite è dato con la cifra che lo mostra. Salvo diversa indicazione le frequenze sono simulazioni a sensore fermo; gli intervalli sono nel capitolo~\ref{sec:val}.
\begin{enumerate}[label=\arabic*.,ref=\arabic*,leftmargin=1.8em]
\item \textbf{Vetro e ostacoli che non rimandano segnale.} Il raggio attraversa la lastra e legge il suolo dietro: 654 punti della lastra su 800 sono dati per liberi a $2{,}0\unita{m}$ e 600 su 800 a $1{,}5\unita{m}$ \prova{S}. Un palo o una lastra scuri senza ritorno risultano «libero» su profondità simulata dalle schede tecniche; la regola dell'impronta regge un corpo esteso senza ritorno (0 fotogrammi con una cella libera su 900 \prova{S}), ma il palo di $3\unita{cm}$ non è stato provato con essa \prova{N}.
\item \textbf{Suolo non visibile vicino al sensore.} Il suolo è visto da $1{,}24\unita{m}$ sull'asse e da $1{,}71\unita{m}$ ai bordi del campo ($1{,}22$ e $1{,}70\unita{m}$ con la focale di 424 pixel, \S\ref{sec:geo-suolo}) \prova{C}. Non sono visti un ostacolo di $15\unita{cm}$ a meno di $1{,}09\unita{m}$ né la testa di una persona di $1{,}70\unita{m}$ a meno di $2{,}01\unita{m}$ \prova{C}; un oggetto alto $0{,}35\unita{m}$ fra $0{,}5$ e $0{,}8\unita{m}$ non riceve alcun pixel e il settore dichiara comunque il suolo letto più lontano \prova{S}.
\item \textbf{Ostacoli sospesi, sottili o più bassi del sensore.} Sotto un ostacolo sospeso il suolo resta «libero» (374 celle nella scena di prova \prova{S}); con una barra di $1\unita{cm}$ a quota $0{,}8\unita{m}$ la distanza dichiarata va oltre la barra nel 4,7--7,3\,\% dei settori per fotogramma \prova{S}. Dietro un ostacolo più basso dell'altezza minima della Tabella~\ref{tab:geo-amin} (circa $1{,}65\unita{m}$ per una statura di $1{,}75\unita{m}$ con il pannello a $1{,}9\unita{m}$) il sistema tratta come nascosta una persona ancora in parte visibile; con un pannello di $1{,}6\unita{m}$ la testa spunta in 10 passaggi su 90 \prova{S}.
\item \textbf{Secondo soggetto nella stessa ombra.} Un secondo soggetto mai visto che esce dalla stessa ombra riceve l'identità del primo: 2400 casi su 2400 nelle prove di base con modello ideale del sensore, fra $0{,}91$ e $0{,}995$ nelle prove d'insieme con modello ideale (3 disposizioni, oltre 3000 casi ciascuna), 19 su 20 (76--99\,\%) su profondità simulata dalle schede tecniche \prova{S}. Nessuna delle regole provate lo ripara: per questo l'identità è «presunta» e l'avviso non dipende da essa.
\item \textbf{Chi svolta, chi esce dal lato lontano, chi parte da fermo.} Per chi svolta la quota dentro $\Reg_\alpha$ nelle tre geometrie è $0{,}86/0{,}67/0{,}79$ a $2\unita{s}$ e $0{,}84/0{,}58/0{,}67$ a $3\unita{s}$; per chi esce dal lato lontano $0{,}82/0{,}61/0{,}79$ e $0{,}76/0{,}57/0{,}70$ \prova{S}, con modello ideale del sensore (\S\ref{sec:val}); con soggetto puntiforme e altri semi, chi esce dal lato lontano scende a $0{,}32$--$0{,}67$ a $2\unita{s}$ e a $0{,}13$--$0{,}47$ a $3\unita{s}$; per chi parte da fermo accanto all'ostacolo, $0{,}11/0{,}16/0{,}26$ a 1, 2 e $3\unita{s}$ nelle prove di base ($N=141$, 134, 54) e $0{,}82$--$0{,}83$ nella prova d'insieme, sugli istanti con il centro nascosto: in questi casi non si dichiara un livello e copre solo l'avviso. Con un ostacolo largo $1\unita{m}$ il livello è «almeno 87 su 100» fino a $1\unita{s}$, poi nessuna cifra.
\item \textbf{Errore del bordo d'ombra.} Con il bordo stimato e non riconosciuto la quota a $1\unita{s}$ scende a $0{,}30$--$0{,}46$ (estremo basso: chi prosegue); a $0{,}74$--$0{,}80$ se il bordo è riconosciuto in 8 cicli su 10 \prova{S}. Con l'ombra stimata più stretta di $0{,}10/0{,}20/0{,}45/0{,}60\unita{m}$ i passaggi in cui il nascosto esce dall'avviso sono $2/15/1\,704/1\,979$ su 2\,000 \prova{S}. Su profondità simulata dalle schede tecniche l'errore del bordo va da $+2$ a $+15\unita{cm}$ al rumore nominale e da $-65$ a $+73\unita{cm}$ al rumore limite.
\item\label{lim:posa} \textbf{Errore di posa.} Un errore d'inclinazione $\delta\theta$ sposta la quota ricostruita di circa $\rho\,\delta\theta$ (\S\ref{sec:geo}): un ostacolo di $10\unita{cm}$ scende sotto la soglia di $7\unita{cm}$ con $0{,}67^\circ$ a $2{,}55\unita{m}$ e $0{,}57^\circ$ a $3\unita{m}$ \prova{C}. Con $+1^\circ$ lo stesso ostacolo a $2{,}55\unita{m}$ risulta «libero» in 7 celle su 7, e una scatola di $10\unita{cm}$ fra $2{,}6$ e $2{,}9\unita{m}$ in 145--150 fotogrammi su 300 (tre esecuzioni) a $-1^\circ$ \prova{S}; le due prove usano convenzioni di segno opposte, e in entrambe l'errore abbassa la quota ricostruita. A $+0{,}5^\circ$ l'identità al ritorno scende a 6 passaggi su 10: la tolleranza è minore di $0{,}5^\circ$.
\item \textbf{Avviso che non si spegne.} Dietro un ostacolo fermo nessuna cella dell'avviso è rilasciabile (0 su 808; 140 passaggi su 140) \prova{S}: l'avviso finisce quando l'ombra si sposta. Nella realizzazione provata, rimosso l'ostacolo, l'avviso residuo resta acceso dopo $2\unita{s}$ in 10 prove su 10 (72--100\,\%), per una causa non determinata: che l'avviso termini quando il suolo è rivisto dopo la rimozione dell'ostacolo non è sostenuto \prova{S} (\S\ref{sec:occ}).
\item \textbf{Profondità realistica.} Su profondità simulata dalle schede tecniche, con la configurazione di riferimento, l'identità al ritorno tiene in 0 passaggi su 90 e il pannello fermo è etichettato mobile in 79 su 90; con i gruppi a celle forti (pannello mobile in 0 passaggi su 140 nel dominio e in 7 su 10 fuori) e l'avviso che resta (regole attive solo dopo le prove d'ingresso, \S\ref{sec:lim-salute}), l'identità tiene in 176 passaggi su 180 [94,4--99,1\,\%] a 424 pixel e in 174 su 180 [92,9--98,5\,\%] a 848, in un dominio stretto (un adulto, $0{,}5$--$1{,}6\unita{m/s}$, entro $2{,}6\unita{m}$) e cade fuori (0--1 su 10 a $0{,}3\unita{m}$ dal pannello; 2 su 10 per un soggetto alto $1{,}0\unita{m}$) \prova{S}. Il simulatore sottostima i pixel senza dato ($7{,}8\,\%$ al livello limite contro $24{,}1\,\%$ e $34{,}4\,\%$ in due registrazioni pubbliche di camere a stereo attivo \prova{P}; \S\ref{sec:val}): le cifre al rumore nominale sono ottimistiche. Sulla stessa profondità simulata la regola della regione persistente (\S\ref{sec:occ}) non è stata eseguita: nella catena provata, che conserva uno stato «perso» senza regione, a soggetto nascosto la traccia è «previsto» in 1\,568 fotogrammi su 4\,552 e «perso» in 2\,983 a 424 pixel di larghezza (1\,584 e 2\,968 su 4\,553 a 848; a entrambe le risoluzioni le due voci sommano a un fotogramma meno del totale); la traccia esiste in tutti i fotogrammi, ma solo 44 e 43 passaggi su 140 hanno una previsione in ogni fotogramma nascosto \prova{S}. La continuità è quindi verificata solo con modello ideale del sensore.
\item \textbf{Classificazione.} Il classificatore su più classi è un componente richiesto dal sistema; in questo lavoro non è realizzato né valutato: le etichette sono assegnate a mano, eventualmente con errori iniettati per misurare la sensibilità. Nessuna classe diversa da «persona» è stata simulata nella catena completa: i percorsi del predittore e tutte le prove riguardano pedoni. Con $3\unita{m}$ di portata e un tempo di reazione ipotizzato $T_r=1{,}0\unita{s}$ l'avviso è tempestivo, per calcolo, fino a $1{,}44\unita{m/s}$: rientrano a velocità tipica 3 classi su 15 e a velocità massima nessuna \prova{C} (\S\ref{sec:occ}). Un'etichetta non produce un «libero» e non toglie né restringe un avviso acceso (invarianti 5 e 15); può cambiare l'esito di un riaggancio, e quindi gli ingressi successivi (\S\ref{sec:occ-classe}).
\item \textbf{Nessuna misura su sensore fisico, nessuna prova con persone.} Rumore, buchi, pixel misti, latenza e costo sulla scheda di bordo non sono misurati; il tempo per ciclo, $6{,}35\unita{ms}$ di media su profondità simulata, sale a $81\unita{ms}$ su una registrazione pubblica affollata \prova{P} (\S\ref{sec:cos}). Lo schermo e i segnali sonori non sono stati provati con persone, vedenti o no; la stessa carenza è documentata nella letteratura sugli ausili, dove su 70 lavori il $51{,}4\,\%$ è valutato in simulazione e il $22{,}9\,\%$ con persone con disabilità visiva \cite{romlay2021}.
\end{enumerate}

\subsection{Limiti a lunga portata}\label{sec:lim-lunga}

Le condizioni di validità e i limiti di questo capitolo riguardano la scala a corto raggio. Per i veicoli veloci valgono le condizioni del \S\ref{sec:lp-validita}, e i limiti che seguono, dati per primi \prova{S}\prova{C}:
\begin{enumerate}[label=L\arabic*.,leftmargin=2.2em]
\item \textbf{Nessun sensore fisico.} Il $46\,\%$ dei parametri dei sensori è un'ipotesi: tutte le probabilità di rilevamento, i falsi allarmi e le latenze; lo stesso vale per i $45\unita{m/s}$ e i $25\unita{m/s^2}$ del drone veloce e per la sezione radar del multirotore.
\item \textbf{Aerei in volo e treni veloci.} L'avviso non è tempestivo (servono $3{,}4\unita{km}$ all'operatore di sito e $7{,}57\unita{km}$ al macchinista, contro $0{,}88\unita{km}$ di «libero» per un multirotore), ed è dichiarato.
% 2C: cifra del «libero» con il pannello.
\item \textbf{«Libero» falso.} Con il radar a pannello, nella configurazione del sito, in 163 passaggi su 171; con una mappa degli ostacoli dichiarata completa ma sbagliata, dietro l'edificio che non vi compare. Il «libero» vale solo per le classi dichiarate.
\item \textbf{Avviso che perde il mezzo.} Con il solo radar per droni l'avviso perde il drone che rallenta o si ferma in 12 passaggi su 200; con il filtro Doppler a bordo in 9 su 100.
% 2C: cifre della collisione.
\item \textbf{Eventi improvvisi.} La collisione è dedotta entro $2\unita{s}$ in 49 passaggi su 60; la zona dei detriti copre il sito e non si riduce; la fiammata è riconosciuta con un innesco ottimista.
\item \textbf{Pochi mezzi reali.} Un solo multirotore e un solo ala fissa nelle traiettorie pubbliche; 43 aerei, 30 auto e 19 treni in prova, in pochi luoghi e giornate. L'inviluppo di $210\unita{m/s}$ non è stato provato di nuovo sulle traiettorie.
\end{enumerate}

\subsection{Estensione a un sensore in moto: analisi dei requisiti}\label{sec:lim-moto}

Con il sensore in moto un residuo fra due fotogrammi può venire dal moto dell'oggetto, dall'errore del moto proprio, dalla parallasse, da un'occlusione che cambia o dal rumore. La regola «mobile» (\S\ref{sec:sti}) confronta due medie di 3 fotogrammi i cui capi distano $T_w=13\Delta=0{,}867\unita{s}$ con la soglia $5\,\sigma\sqrt{2/3}$: uno spostamento apparente $\delta$ dovuto alla posa non compensata produce un'etichetta falsa con probabilità $Q(5-\delta/(\sigma\sqrt{2/3}))+Q(5+\delta/(\sigma\sqrt{2/3}))$, che vale $10\,\%$ per $\delta=3{,}036\,\sigma$ e $1\,\%$ per $\delta=2{,}183\,\sigma$ \prova{D} (\S\ref{sec:sti-mobile}). Ne seguono i requisiti della tabella~\ref{tab:lim-moto}, che renderebbero valide le stesse regole; sono calcoli in forma chiusa nel modello a errori gaussiani indipendenti fra fotogrammi \prova{C}. Il solo controllo numerico è un Monte Carlo della formula nello stesso modello d'errore ($N=2\cdot10^6$ valutazioni: $0{,}0999$ e $0{,}0100$). Nessun requisito è verificato su una piattaforma in moto né con una simulazione del sensore in moto, e gli errori reali delle sorgenti di posa non sono noti \prova{N}: la tabella è un'analisi, non un risultato.

\begin{table}[H]\centering\small
\caption{Requisiti sulla posa per un sensore in moto (soglia: $10\,\%$ di falsi, salvo dove indicato) \prova{C}.}\label{tab:lim-moto}
\begin{tabularx}{\linewidth}{@{}Lp{3.4cm}p{3.4cm}@{}}
\toprule
Requisito & a $2\unita{m}$ & a $3\unita{m}$\\
\midrule
Imbardata relativa sulla finestra $T_w$ & \multicolumn{2}{@{}l}{$0{,}88^\circ$ ($1{,}01^\circ/\mathrm{s}$) a ogni distanza; $0{,}63^\circ$ ($0{,}73^\circ/\mathrm{s}$) per l'$1\,\%$}\\
Errore di velocità propria, trasversale / radiale & $3{,}5$ / $5{,}3\unita{cm/s}$ & $5{,}3$ / $11{,}9\unita{cm/s}$\\
Ritardo fra posa e profondità & \multicolumn{2}{@{}l}{$29\unita{ms}$ con un gradino di $30^\circ/\mathrm{s}$; $74\unita{ms}$ con un'oscillazione di $\pm3^\circ$ a $1\unita{Hz}$}\\
Inclinazione rispetto al suolo ($10\,\%$ di celle di suolo oltre la soglia di quota) & $1{,}70^\circ$; $1{,}27^\circ$ con $\pm2{,}5\unita{cm}$ d'oscillazione verticale & $1{,}03^\circ$; $0{,}79^\circ$\\
Scostamento non modellato perché la regione copra ancora $0{,}90$ & \multicolumn{2}{@{}l}{$9{,}7\unita{cm}$ a $0{,}67\unita{s}$; $16\unita{cm}$ a $1\unita{s}$; $41\unita{cm}$ a $2\unita{s}$}\\
Quota di fascia lasciata alla posa ($0{,}10\unita{m}$) & $2{,}9^\circ$; ritardo $95\unita{ms}$ a $30^\circ/\mathrm{s}$ & $1{,}9^\circ$; $64\unita{ms}$\\
\bottomrule
\end{tabularx}
\end{table}

\paragraph{Che cosa resta invariato e che cosa no.} La regola «mai libero dove non è osservato» è calcolata sul singolo fotogramma: non usa la posa relativa fra fotogrammi e resta valida in moto. Lo stesso vale per il divieto di produrre un «libero» da memoria, etichette o timer, e per la marca «osservato» o «dedotto». Dipendono invece dalla posa: la correttezza di ciò che è dichiarato «libero», che richiede fotogramma per fotogramma l'inclinazione rispetto al suolo entro la tolleranza del limite~\ref{lim:posa}; la regola «mobile», che senza compensazione del moto proprio scambia il moto apparente per moto dell'oggetto (nel seguimento di più oggetti la compensazione del moto della camera è usata per questo \cite{aharon2022}); l'ombra, che non è più quella di un ostacolo fermo; la fascia di guardia, che cresce con $v_e$. Una memoria dello spazio osservato tornerebbe necessaria, perché i fotogrammi vedono raggi diversi, e la forma compatibile con gli invarianti è una griglia locale che segue il sensore \cite{usenko2017}, usata a senso unico: può aggiungere un «occupato ricordato», non può produrre un «libero» né sciogliere una possibile presenza (\S\ref{sec:cos}); la conferma su tre fotogrammi vale solo sotto $0{,}375\unita{m/s}$ di traslazione propria \prova{C}.

\subsection{Misure che deciderebbero i punti non verificati sperimentalmente}\label{sec:lim-misure}

La tabella~\ref{tab:lim-misure} descrive, come metodo di verifica, le misure su un sensore fisico che deciderebbero gli enunciati marcati \prova{N}. Le soglie sono fissate prima della misura; dove le fonti non fissano una numerosità la casella lo dice.

\begin{table}[H]\centering\footnotesize
\caption{Protocollo di verifica: grandezza, numerosità, soglia di smentita.}\label{tab:lim-misure}
\begin{tabularx}{\linewidth}{@{}rLp{4.0cm}p{4.6cm}@{}}
\toprule
 & Grandezza misurata & Numerosità & Soglia\\
\midrule
1 & Gruppi e tracce sul suolo vuoto, a posa tarata; stato di un oggetto di $10\unita{cm}$ a $2{,}5\unita{m}$ & $30\unita{s}$ & 0 gruppi, 0 tracce confermate; oggetto «occupato»\\
2 & Celle «libero» dentro l'ombra vera di un pannello, tracciata sul suolo & 150 fotogrammi & 0 celle\\
3 & Quota dei pixel dell'impronta entro il margine di lettura del suolo, a scena ferma & $60\unita{s}$ & sotto l'$80\,\%$ oltre $2\unita{m}$: tarare il margine prima di dichiarare una distanza\\
4 & Quota dei cicli in cui il bordo dell'ombra è riconosciuto; errore del bordo sui due lati del pannello & non fissata & condizione per dichiarare un livello: errore entro $\pm20\unita{cm}$\\
5 & Persona dentro $\Reg_\alpha$ a 1 e $2\unita{s}$, con occlusione virtuale (una copia della catena riceve le misure con un cuneo cancellato; la persona resta visibile e fa da verità) & 100 passaggi nel dominio dichiarato, camminatori diversi da quelli della taratura & almeno 87 su 100: «più di 8 su 10» (esclude $0{,}80$ con potenza $0{,}88$ se il vero è $0{,}90$); almeno 96 su 100: «almeno 9 su 10»\\
6 & Segni pieni a soggetto nascosto, nella stessa prova & gli stessi 100 passaggi & 0\\
7 & Avvisi rilasciati senza aver rivisto il suolo, con un secondo soggetto; avviso dopo la rimozione dell'ostacolo & non fissata & 0 rilasci; avviso spento entro 5 fotogrammi e 0 nascosti fuori dall'avviso\\
8 & Latenza dall'evento all'uscita sonora & non fissata & 95° percentile entro $150\unita{ms}$\\
9 & Segnale di affidabilità persa coprendo il sensore per 1 e $2\unita{s}$; avvio con i filtri accesi o con la posa spostata di $1^\circ$ & $60\unita{s}$ per i falsi segnali & segnale entro $1\unita{s}$; 0 falsi; il sistema non parte\\
10 & Uscite di due esecuzioni sugli stessi fotogrammi che differiscono solo nelle etichette di classe (permutate, assenti, cambiate in occlusione) & non fissata & celle «libero» e distanze identiche; per ogni ingresso riconosciuto in entrambe le esecuzioni, avviso con etichette non più stretto; nessun livello più basso\\
\bottomrule
\end{tabularx}
\end{table}

Una prova con dieci passaggi non distingue $0{,}80$ da $0{,}90$ (\S\ref{sec:val}); trenta passaggi senza errori non bastano a dichiarare «9 su 10» (limite inferiore $0{,}886$). Il tempo di reazione di una persona con disabilità visiva a un avviso sonoro, da cui dipende la riga corrispondente della tabella~\ref{tab:lim-portatore}, richiede una misura con persone, fuori da questo protocollo.
